\chapter{后续实施计划与阶段验收}\label{chap:middle-plan}
\section{以基础平台恢复建立实物验证入口}

下一阶段首先形成原 AUV 当前状态的可核验记录，覆盖能源、通信、必要反馈与执行接口。现有历史设计资料用于识别预期连接关系，现场检查则确定当前实际配置，两者存在差异时以实测和版本记录更新实验条件。重点不是将全部设备一次接入，而是明确受控运行所需的最小完整链路，并确保每次增加功能都能够退回已经验证的状态。

能源检查应为后续装配和通电提供准入依据，通信检查应确认命令、反馈与异常状态的方向和含义，执行链检查应形成通道极性、响应和停机条件记录。必要反馈、机械固定和水密不能因任务载荷后置而省略。完成这些前置条件后，再由具有现场条件和授权的人员安排受控水域试验，验证基础运动、退出与恢复处置。具体操作参数由现场方案确定，本报告不将软件生成的设定值视为执行许可。

基础阶段以 M2.5 所对应的“基础平台恢复与水中受控运行验证”为里程碑。其验收需要同次记录中同时出现有效反馈、受控命令、实际运动和退出结果；仅有上电照片、通信连通或台架动作仍属于前置进展。此阶段主要回答平台是否已具备承载后续实验的条件，海缆自主巡检的完成另在任务阶段评价。

\section{磁载荷集成与完整任务验证}

平台具备稳定运行条件后，磁载荷集成按静态、受控运动和任务闭环逐级推进。首先固定传感器位置、外参、供电及布线方案，记录计算设备和执行机构在不同工作状态下的磁背景。随后在具有可溯源相对几何的等效电缆条件下进行重复横切，建立频带、姿态、扫描范围与宽度估计之间的关系。实验应保留无目标、邻近干扰及不满足场形条件的对照，使有效率与误差能够共同评价。

在线阶段先验证采集、时间戳、估计与记录是否在同一时间轴内一致，再逐步接入运动决策。若利用声学或调查信息限定目标，应同时记录目标关联依据及其覆盖范围。未获得深度条件时继续保留空值，不用后处理真值补填在线结果。这样可把横向巡检验证与埋深标定分开推进，同时保持输出含义一致。

完整任务应覆盖搜索、确认、接入、跟踪及异常后的处置，并以完整路线而非选定稳定片段计算误差。记录中需同时包含任务输入、导航与磁观测、算法状态、实际命令和运动反馈。故障试验除了核查命令撤回，还要关注载体在撤回后的运动和边界距离，从而确定下一阶段需要的停止、回收或其他恢复措施。

\section{工作包、依赖和完成条件}

表\ref{tab:middle-workpackages}给出实施顺序和建议窗口。表中时间采用前置条件满足后的相对周数，是本次组织的资源规划建议；实际进入时间取决于设备、人员和场地条件，不构成海试或答辩日期承诺。磁标定资料整理与运行日志工具可在基础恢复阶段并行准备，实物接管和水中任务则必须等待相应准入完成。

\begingroup\small
\begin{longtable}{p{2.3cm}p{11.4cm}}
\caption{后续工作包及逐级验收安排}\label{tab:middle-workpackages}\\
\toprule 工作包&当前基础、实施条件和验收\\\midrule
\endfirsthead
\multicolumn{2}{c}{表\thetable\quad 后续工作包及逐级验收安排（续）}\\
\toprule 工作包&当前基础、实施条件和验收\\\midrule
\endhead
\bottomrule\endfoot
W1 基础状态与能源通信&
\textbf{当前基础：}历史平台资料、真实通信与零执行器安全子链记录。
\textbf{待完成：}现场配置核对，能源、必要反馈及接口状态确认。
\textbf{依赖：}可接触平台、设备资料及现场授权。
\textbf{窗口：}具备现场条件后的第1阶段，建议1--2周。
\textbf{验收：}版本和接线记录、能源检查、命令反馈与异常状态日志。
\textbf{阻塞处理：}缺失部件或资料单独登记，先完成不依赖该项的接口核查，不跨越通电准入。\\[5pt]
W2 执行与水中受控运行&
\textbf{当前基础：}执行脑接口、控制及故障状态传播方法。
\textbf{待完成：}极性与响应确认、装配水密、基础运动和退出试验。
\textbf{依赖：}W1、机械固定、水密、必要反馈及受控试验条件。
\textbf{窗口：}W1通过后的第2阶段，建议1--2周。
\textbf{验收：}同次命令、运动、反馈和退出记录，按现场方案核查M2.5。
\textbf{阻塞处理：}继续台架或影子验证，水中受控运行保持未完成。\\[5pt]
W3 磁载荷与几何标定&
\textbf{当前基础：}TMR/ADC采集和实验室相对导线拟合。
\textbf{待完成：}装艇外参、磁背景、可溯源横切及重复性标定。
\textbf{依赖：}固定安装、稳定采集、已知几何和可核验激励；运动采集依赖W2。
\textbf{窗口：}台架准备可与W1并行，装艇后建议1--2周。
\textbf{验收：}原始波形、参考几何、背景对照、误差及拒绝样本。
\textbf{阻塞处理：}保留可检验的相对量；深度或目标身份无依据时不发布。\\[5pt]
W4 在线任务与安全退出&
\textbf{当前基础：}限定条件仿真及ROS接口与独立记录方法。
\textbf{待完成：}目标平台完整时序、在线搜索横切跟踪、异常处置。
\textbf{依赖：}W2、W3及目标参考、导航条件、场地和操作方案。
\textbf{窗口：}前置通过后的第3阶段，建议2--3周。
\textbf{验收：}完整路线误差、阶段转换、同次输入到执行记录及撤回后运动。
\textbf{阻塞处理：}隔离到采集、影子或单回路验证，记录实际达到的层级。\\[5pt]
W5 重复对照与论文收口&
\textbf{当前基础：}冻结仿真、实物子链及本报告证据索引。
\textbf{待完成：}重复试验、误差归因、失败汇总和最终章节更新。
\textbf{依赖：}前述实验版本稳定、原始数据与配置齐备。
\textbf{窗口：}随W3/W4滚动整理，最终数据形成后建议1--2周。
\textbf{验收：}图表、统计分母、条件和结论一致，证据可回查。
\textbf{阻塞处理：}按截止日冻结已达成层级，未完成任务及影响单列，不改变原选题要求。\\
\end{longtable}
\endgroup

\section{阶段更新与预期成果}

各阶段统一保存设备与软件版本、配置、原始数据、处理方法和结果说明。每次运行明确正常、受限或中止状态，保留失败样本和退出原因；更换硬件、安装位置或控制参数后，记录需要重新检验的条件。这样，已有结果能够为下一步提供依据，同时避免把不同版本下的局部最佳表现组合成完整性能。

中期之后的新证据优先更新平台进展、共性问题和工作包状态。若基础恢复先于载荷集成完成，应首先加入受控运动及退出的同次记录；若磁标定先完成，则更新几何模型适用范围与误差预算。新增结果只有在明确验证层级后才进入总结，未达到的整艇或海缆任务继续列为待完成。替代试验可以支撑对应子问题，其作用和局限应与所替代的条件一并说明。

预期通过上述顺序，将目前分层形成的方法和子链接口逐步连接到同一实物平台。最终研究成果应能够说明：在给定观测和运动条件下可以估计什么、何时能够执行巡检、何时应中止或转交，以及这些判断由哪些数据支持。完整任务仍以实际实验为依据；本报告形成的结构、冻结证据与验收路线为后续推进提供可更新的基础。
